\section{Discussion}

\subsection{Evidential Methods for Neural Operators}

Our experiments establish evidential deep learning as the preferred approach for uncertainty quantification in 
neural operators. Key advantages include:

\textbf{Superior calibration}: Evidential methods (ECE $\sim$ 0.004) outperform MC Dropout (0.019) and Deep 
Ensembles (0.011) while maintaining 93--95\% coverage. This superiority holds across both Navier-Stokes and 
Darcy Flow, with the performance gap widening on harder problems.

\textbf{Computational efficiency}: Single-pass inference provides 5--30$\times$ speedup over alternatives, 
critical for real-time applications and high-throughput scientific computing workflows.

\textbf{Informative uncertainty}: Uncertainty-error correlation of 0.61 (vs. 0.31--0.33 for baselines) indicates 
evidential uncertainty better reflects prediction quality.

\textbf{MC Dropout unsuitable}: Across all metrics, MC Dropout exhibits inferior performance 
(6$\times$ worse MSE, lowest correlation), confirming it should be avoided for neural operators. 
The spatial coherence of PDE solutions may be incompatible with dropout's stochastic regularization.

\subsection{The Calibration-Decomposition Trade-off}

A fundamental limitation emerges from our decomposition analysis: methods achieving tight calibration 
(ECE $\sim$ 0.004) exhibit invalid decomposition (aleatoric variation $>$ 89\%), while valid decomposition 
(KL-Divergence) yields poor calibration (ECE = 0.030).

This trade-off appears intrinsic to the NIG parameterization. During training, the network increases evidence ($\nu, \alpha$) for correct predictions, directly reducing epistemic uncertainty ($\propto 1/\nu$). However, aleatoric uncertainty ($= \beta/(\alpha-1)$) also decreases because the loss function optimizes $\beta$ to match observed prediction variance. As the model improves with more data, prediction errors decrease, so the optimal $\beta$ decreases. Since $\alpha$ increases while $\beta$ decreases, the ratio $\beta/(\alpha-1)$ drops substantially---even though it theoretically should remain constant as irreducible data noise. The NIG parameterization lacks separate constraints to preserve aleatoric stability.

\textbf{Practical implication}: Practitioners should treat evidential outputs as \emph{total uncertainty} 
rather than interpreting the decomposition. For applications requiring true epistemic-aleatoric separation, 
Deep Ensembles remain preferable despite computational cost.

\subsection{Reinterpreting OOD Detection}

The stark contrast between parameter-level OOD (AUROC = 0.55) and operator-level OOD (AUROC = 1.0) suggests a 
domain-specific interpretation: parameter extrapolation within the same PDE family may not constitute true 
distributional shift from the operator learning perspective.

The FNO learns the Navier-Stokes operator in Fourier space, which generalizes across Reynolds numbers despite 
different flow characteristics. This explains maintained confidence at Re = 5000. Conversely, cross-PDE transfer 
represents genuine operator-level shift where the governing equations differ fundamentally, yielding appropriate 
massive uncertainty increases (36--378$\times$).

This finding has deployment implications: evidential uncertainty reliably monitors operator-level distributional 
shift but should not be solely relied upon for parameter extrapolation validation.

\subsection{Architectural Insights}

Our ablation study identifies critical components for evidential FNOs:
\begin{itemize}
    \item \textbf{Essential}: Skip connections and GELU activations---removal causes catastrophic failure
    \item \textbf{Beneficial}: Global residual connections (+5\% improvement)
    \item \textbf{Detrimental}: Batch normalization (2.2$\times$ worse), likely interfering with NIG parameter 
    learning
    \item \textbf{Flexible}: Evidential head complexity matters little---shallow heads maintain quality with fewer 
    parameters
\end{itemize}

\subsection{Limitations}

This study has several limitations: (1) evaluation on two PDEs (Navier-Stokes, Darcy Flow) may not generalize to 
other PDE families; (2) single architecture (FNO) tested---other operators may behave differently; (3) 
synthetic training data may not capture real-world measurement noise; (4) limited OOD scenarios tested. 
Future work should validate findings across broader PDE classes, architectures, and real-world applications.

\section{Conclusion}

This paper presents comprehensive comparison of uncertainty quantification methods for Fourier Neural Operators,
 evaluating MC Dropout, Deep Ensembles, and Evidential Deep Learning across calibration, efficiency, 
 decomposition validity, and OOD detection.

\textbf{Key findings}:
\begin{enumerate}
    \item Evidential methods achieve state-of-the-art calibration (ECE = 0.004) with single-pass inference, 
    outperforming MC Dropout (5$\times$ better ECE) and Deep Ensembles (3$\times$ better ECE) while being 
    5--30$\times$ faster.

    \item Uncertainty decomposition fails for well-calibrated evidential methods---both epistemic and aleatoric 
    uncertainties decrease with training data, revealing a fundamental calibration-decomposition trade-off in 
    the NIG parameterization.

    \item OOD detection capability depends critically on shift type: near-random for parameter extrapolation 
    (AUROC = 0.55) but perfect for operator-level shift (AUROC = 1.0), suggesting neural operators learn 
    physics-informed representations.

    \item MC Dropout should be avoided for neural operators due to consistently poor performance across all metrics.
\end{enumerate}

\textbf{Recommendations}: For production deployment, use evidential FNO with Uncertainty-Aware or Improved 
regularization. Report total uncertainty only (ignore decomposition). Implement threshold-based fallback for 
operator-level OOD detection. Reserve Deep Ensembles for applications requiring robust epistemic uncertainty despite 
computational cost.

Future work should investigate alternative parameterizations overcoming the calibration-decomposition trade-off, 
extend evaluation to broader PDE families and architectures, and develop physics-informed OOD metrics for parameter 
extrapolation scenarios.